\documentclass[11pt,a4paper]{article}

\usepackage[margin=2.3cm]{geometry}
\usepackage{amsmath,amssymb}
\usepackage{booktabs}
\usepackage{longtable}
\usepackage{array}
\usepackage{xcolor}
\usepackage[most]{tcolorbox}
\usepackage{fancyhdr}
\usepackage{titlesec}
\usepackage{enumitem}
\usepackage{listings}
\usepackage{hyperref}

\definecolor{strivenavy}{HTML}{1A3D6D}
\definecolor{striveteal}{HTML}{2E8B8B}
\definecolor{strivewarn}{HTML}{C0392B}
\definecolor{strivegold}{HTML}{B8860B}
\definecolor{strivegrey}{HTML}{F4F6F8}

\hypersetup{colorlinks=true,linkcolor=strivenavy,urlcolor=striveteal,
  pdftitle={STRIVE: Detection Handoff},
  pdfauthor={Team STRIVE, IISER Bhopal}}

\titleformat{\section}{\Large\bfseries\color{strivenavy}}{\thesection}{0.7em}{}
\titleformat{\subsection}{\large\bfseries\color{striveteal}}{\thesubsection}{0.6em}{}

\pagestyle{fancy}
\fancyhf{}
\lhead{\small\color{strivenavy}STRIVE --- Detection Handoff}
\rhead{\small\color{strivenavy}SIH26104}
\cfoot{\thepage}
\renewcommand{\headrulewidth}{0.4pt}

\newtcolorbox{donebox}[1]{colback=green!3,colframe=striveteal,
  fonttitle=\bfseries,title={#1},boxrule=0.6pt,arc=2pt,left=4pt,right=4pt}
\newtcolorbox{todobox}[1]{colback=yellow!8,colframe=strivegold,
  fonttitle=\bfseries,title={#1},boxrule=0.6pt,arc=2pt,left=4pt,right=4pt}
\newtcolorbox{warnbox}[1]{colback=red!3,colframe=strivewarn,
  fonttitle=\bfseries,title={#1},boxrule=0.6pt,arc=2pt,left=4pt,right=4pt}

\lstset{basicstyle=\ttfamily\footnotesize,frame=single,framesep=4pt,
  backgroundcolor=\color{strivegrey},rulecolor=\color{gray},breaklines=true,
  showstringspaces=false}

\newcommand{\code}[1]{\texttt{\small #1}}

\begin{document}

\begin{titlepage}
\centering
\vspace*{1.8cm}
{\Huge\bfseries\color{strivenavy} STRIVE\par}
\vspace{0.5cm}
{\LARGE\color{striveteal} Detection Handoff\par}
\vspace{0.4cm}
{\large What the platform provides, and how to plug a real detector into it\par}
\vspace{2cm}
\begin{tabular}{rl}
\textbf{Problem Statement} & SIH26104\\
\textbf{Organization} & AICTE Cyber Security Cell\\
\textbf{Theme} & Blockchain \& Cybersecurity\\
\textbf{Team} & STRIVE, IISER Bhopal\\
\textbf{Audience} & Detection / ML teammates\\
\end{tabular}
\vfill
\begin{tcolorbox}[colback=strivegrey,colframe=strivenavy,width=0.92\textwidth]
\textbf{One-line summary.} The streaming platform, channel intelligence, codec
robustness lab and evaluation harness are built and measured. The detector is a
DSP placeholder scoring \textbf{AUC 0.708} on real ASVspoof data. Everything you
need to replace it is described in \S\ref{sec:interface}; the corpus and GPU are
already provisioned.
\end{tcolorbox}
\vspace{1cm}
\end{titlepage}

\tableofcontents
\newpage

%==========================================================
\section{Where the project stands}

\subsection{Built and measured}

\begin{donebox}{Streaming and real-time}
Browser microphone $\to$ canonical 16\,kHz mono $\to$ 20\,ms binary PCM
(\code{pcm-v2}) $\to$ independent server ingestion and scoring. Configurable
window/hop geometry (2\,s/1\,s, 2\,s/0.5\,s, 4\,s/0.5\,s all tested). Multi-rate
branch scheduler. Bounded capture queue so inference stalls cannot block audio.
Optional Rust audio runtime. Per-stage latency instrumentation throughout.
\end{donebox}

\begin{donebox}{Channel intelligence}
Per-window occupied bandwidth, SNR proxy, clipping ratio, RMS dBFS, cross-window
discontinuity and a conservative reliability scalar. Verified to separate G.711,
Opus and narrowband transport from clean audio.
\end{donebox}

\begin{donebox}{Codec robustness lab}
Fifteen registered conditions (Opus 6/12/16/24/32\,kbps, G.711 $\mu$-law and
A-law, 8\,kHz narrowband, noise at 20/10/5\,dB, packet loss 1/3/5\%). Paired
manifests that keep \code{source\_id} across every derived clip so no split
leakage is possible.
\end{donebox}

\begin{donebox}{Evaluation harness}
Streaming WAV simulator (deterministic and wall-clock modes), latency benchmark,
and a codec benchmark producing EER, ROC-AUC, TPR at fixed FPR, Brier, score
distributions and per-condition latency. \textbf{183 tests pass, 1 skipped.}
\end{donebox}

\begin{warnbox}{Not built: the detector}
The default extractor is \code{dsp-surrogate-v1}: hand-designed spectral features,
no learned weights, explicitly flagged \code{demo\_only: true} on every event.
Research adapters for the frozen NII countermeasure exist but \textbf{have never
been run} --- no model bundle is installed. Every detection number in this
document comes from the placeholder.
\end{warnbox}

\subsection{Measured detection floor}

ASVspoof 2019 LA, \textbf{unseen-generator} evaluation: reference index built on
attacks A01--A06, evaluated on A07--A19. 20 vs 67 speakers, disjoint. 3{,}000
evaluation clips per condition, $n_0 = 1451$ genuine, zero failures across 18{,}000
evaluations.

\begin{center}
\begin{tabular}{@{}lrrrrr@{}}
\toprule
\textbf{Condition} & \textbf{EER \%} & \textbf{$\Delta$EER} & \textbf{AUC} & \textbf{$\Delta$AUC} & \textbf{TPR@1\%FPR}\\
\midrule
\texttt{clean}          & 35.04 & ---    & 0.708 & ---     & 0.056\\
\texttt{opus\_32k}      & 36.84 & $+1.80$ & 0.672 & $-0.036$ & 0.070\\
\texttt{opus\_16k}      & 37.44 & $+2.40$ & 0.676 & $-0.031$ & 0.046\\
\texttt{g711\_ulaw}     & 37.86 & $+2.82$ & 0.669 & $-0.039$ & 0.075\\
\texttt{narrowband\_8k} & 38.07 & $+3.03$ & 0.670 & $-0.038$ & 0.072\\
\textbf{\texttt{noise\_10db}} & \textbf{48.43} & $\mathbf{+13.39}$ & \textbf{0.533} & $\mathbf{-0.174}$ & 0.103\\
\bottomrule
\end{tabular}
\end{center}

\textbf{Read this as a floor, not an achievement.} AUC 0.708 is weak but
non-trivial signal from unsupervised DSP features. It is the number a trained
countermeasure must beat, and it is published precisely so that any improvement
is measurable rather than asserted.

\subsection{Measured latency}

Live loopback WebSocket, 400 packets, 2\,s window / 0.5\,s hop, i5-1235U CPU.

\begin{center}
\begin{tabular}{@{}lrr@{}}
\toprule
\textbf{Measurement} & \textbf{Normal} & \textbf{1200\,ms injected inference}\\
\midrule
Packets acknowledged & 400/400 & 400/400\\
Ingestion p95 & 0.615\,ms & 0.612\,ms\\
Window-end $\to$ client p95 & 82.53\,ms & 1720.98\,ms\\
Windows scored / discarded & 13 / 0 & 6 / 7\\
Maximum pending depth & 1 & 1\\
\bottomrule
\end{tabular}
\end{center}

\begin{donebox}{What the overload column buys you}
With inference deliberately stalled at 1.2\,s per window, ingestion stays at
0.612\,ms and every packet is still acknowledged. Stale windows are discarded, not
queued, and pending depth never exceeds one.

\textbf{Practical consequence for you: a slow model degrades the update rate, it
does not break the call.} You are free to try a heavy model without first
optimising it. Compute, queue, audio-lag and dropped-window counters make the cost
visible in every event.
\end{donebox}

%==========================================================
\section{The detector interface}
\label{sec:interface}

This is the whole contract. Anything satisfying it drops into the running system
with no other changes.

\subsection{What you must provide}

\begin{lstlisting}[language=Python]
class YourExtractor:
    id: str                 # e.g. "nii-mms300m-<hash12>"; appears in every event
    is_surrogate: bool      # False for a real model -> removes the demo_only flag

    def extract(self, waveform: np.ndarray) -> Features:
        """Exactly window_s * 16000 float32 samples, normalised to [-1, 1]."""

    # Optional
    pipeline_version: str                      # provenance beyond the model hash
    def identify_language(self, audio) -> tuple[str, str]:   # (tag, source)
    def close(self) -> None                                  # release GPU memory
\end{lstlisting}

\subsection{What you must return}

\begin{lstlisting}[language=Python]
@dataclass
class Features:
    cm: np.ndarray            # countermeasure embedding, unit-normalised
    profile: np.ndarray       # speaker/acoustic profile vector, own dimension
    segments: list[Segment]   # time-aligned sub-window vectors
    cm_score: float | None    # direct spoof probability in [0,1], or None
    metadata: dict            # free-form; surfaced in the event's features field

@dataclass
class Segment:
    start_s: float            # relative to window start
    end_s: float
    vector: np.ndarray        # all segments in a window must share a dimension
    token: int = -1           # phoneme/cluster id, -1 if unavailable
\end{lstlisting}

\begin{center}
\begin{tabular}{@{}lp{10.4cm}@{}}
\toprule
\textbf{Field} & \textbf{Consumed by}\\
\midrule
\code{cm} & Track 1, FAISS retrieval against the labelled reference index. Must match the index's dimension and the extractor id it was built with.\\
\code{profile} & Track 2, in-call session consistency. Dimension may differ from \code{cm} and must stay fixed within a call.\\
\code{segments} & Track 3, boundary coherence across the window seam.\\
\code{cm\_score} & Recorded as \code{cm\_head\_score}; not yet fused. See \S\ref{sec:next}.\\
\bottomrule
\end{tabular}
\end{center}

\begin{warnbox}{Three contracts that will fail loudly if broken}
\begin{enumerate}[nosep,leftmargin=1.4em]
\item \textbf{Input shape is exact.} \code{extract} receives precisely
\code{window\_s} $\times$ 16000 samples. Not approximately.
\item \textbf{Dimensions are fixed within a call.} Changing \code{profile} or
segment dimension mid-call raises; the session store cannot mix vector spaces.
\item \textbf{The index must match the extractor.} Startup refuses if
\code{index.metadata["extractor\_id"] != extractor.id}, and research mode refuses a
\code{demo\_only} fixture index. This is deliberate: a silently mismatched index
produces confident nonsense.
\end{enumerate}
\end{warnbox}

\subsection{Where it is called}

\begin{lstlisting}[language=Python]
# strive/engine.py
features = self.extractor.extract(window.samples)        # per window
self.language, src = self.extractor.identify_language(...)  # once, after bootstrap
\end{lstlisting}

Failures are caught and reported as \code{MODEL\_OR\_INDEX\_ERROR} with explicit
uncertainty --- never as a zero score. Set \code{raise\_model\_errors=True} in
config while debugging.

%==========================================================
\section{Findings that should shape your design}

These come from the surrogate, but the mechanisms are model-independent.

\subsection{Noise destroys detection; codecs merely dent it}

Every codec leg costs $0.031$--$0.039$ AUC. Additive noise at 10\,dB SNR costs
$\mathbf{0.174}$ --- about five times worse --- driving AUC to $0.533$ and
collapsing the bonafide/spoof median gap to \textbf{exactly zero}.

\textbf{Implication:} noise robustness deserves more of your effort than codec
robustness. Test under additive noise early; a model tuned only on clean or
codec-degraded audio will look far better than it is.

\subsection{The evidence gate is inverted under noise}

Clips clearing the live-call evidence gate: 114 clean, \textbf{651} at 10\,dB SNR
--- a $5.7\times$ increase. The cause is the demo-mode energy VAD
(\code{strive/audio.py::activity}, RMS $> 0.008$), whose voiced fraction rises from
$0.564$ to $1.000$ under noise. Every frame is counted as speech.

\begin{warnbox}{The system gets more confident exactly when it should get less}
This is safety-relevant and currently unfixed. Research mode substitutes a trained
WebRTC VAD, which should behave better --- \textbf{but that has not been measured}.
Verifying it is a concrete, valuable early task.
\end{warnbox}

\subsection{The surrogate's signal is generator-agnostic}

Seen-generator AUC was $0.707$; unseen-generator is $0.708$. No drop. The
placeholder is not memorising attack families --- it is picking up a generic
synthesis cue.

\textbf{Implication:} a trained countermeasure will probably show a real
seen/unseen gap, because it \emph{can} overfit to known attacks. Always report
unseen-generator numbers. The corpus splits make this automatic: A01--A06 for the
index, A07--A19 for evaluation.

%==========================================================
\section{How to move forward}
\label{sec:next}

\subsection{Already provisioned}

\begin{center}
\begin{tabular}{@{}lp{9.6cm}@{}}
\toprule
\textbf{Resource} & \textbf{Status}\\
\midrule
GPU & \code{stork} (172.28.79.149): RTX 3090, 24\,GB, 61\,GB RAM, 32 threads\\
Corpus & ASVspoof 2019 LA complete, 122{,}299 files, md5-verified, at\newline
\code{/media/storage/moonlab/strive-data/asv2019}\\
Manifests & Official speaker-disjoint splits already built:\newline
reference 25{,}380 / validation 24{,}844 / test 71{,}237\\
Codec matrix & 18{,}000 degraded clips across 6 conditions, ready to reuse\\
Python & 3.12 venv for the app; system 3.10 for the Fairseq export\\
\bottomrule
\end{tabular}
\end{center}

The Python 3.10 on \code{stork} is not an inconvenience --- the NII export needs
legacy Fairseq, which will not install alongside the 3.12 application runtime.

\subsection{Suggested order}

\textbf{Step 1 --- Get one real number.} Read \code{docs/MODELS.md}, then:

\begin{lstlisting}[language=bash]
python scripts/prepare_models.py --download --variant mms-300m --include-language
# separate Python 3.10 environment for the Fairseq export
python scripts/export_nii.py --variant mms-300m \
       --weights models/nii-source/model.safetensors --output models/cm.pt
# seal the bundle: manifest.json pinning sha256 for every file
python scripts/build_index.py data/manifests/asv2019/reference.csv --models models
python scripts/start.py --research
\end{lstlisting}

Start with \textbf{MMS-300M}, not XLS-R-2B. It is 1024-dimensional and far
lighter; get the pipeline working before scaling the model.

\textbf{Step 2 --- Re-run the benchmark and compare against the floor.}

\begin{lstlisting}[language=bash]
python scripts/codec_benchmark.py \
  --matrix data/manifests/asv2019/matrix_full.csv \
  --reference-manifest data/manifests/asv2019/ref_sample.csv \
  --output evidence/codec-benchmark-research.json
\end{lstlisting}

Identical cohort, identical conditions. The delta against AUC $0.708$ is the
contribution, and it is directly quotable.

\textbf{Step 3 --- Fuse \code{cm\_score}.} The research extractor returns a direct
spoof probability that is currently \emph{recorded but unused}; Track 1 relies on
retrieval alone. Wiring it in is likely the single largest accuracy gain available
and needs no new model.

\textbf{Step 4 --- Channel-aware fusion (FUS-03).} Under \code{noise\_10db} the
channel layer reports \code{snr\_db} $= 19.6$ and \code{quality} $= 0.898$ while AUC
collapses to $0.533$. The signal needed to discount that evidence is already
computed and already in every event --- the fusion simply ignores it. Constants
belong in config; the feature needs an on/off flag.

\textbf{Step 5 --- Prosody (BR-04) and identity (BR-05) branches.} The
\code{BranchResult} schema and multi-rate scheduler already exist. Both are
optional by construction: fusion masks unavailable branches and renormalises.

\subsection{Rules we have been holding to}

\begin{enumerate}[leftmargin=1.6em]
\item \textbf{Never tune a threshold to make a demo look better.} Two known
weaknesses (EMA warm-up, mature-call dilution) are pinned by tests on purpose.
\item \textbf{Absent evidence is masked, never zeroed.} A missing or stale branch
must not contribute a score that reads as ``genuine''.
\item \textbf{Unavailable measurements are \code{null}}, never $0$.
\item \textbf{No fabricated numbers.} If a cohort is too small, the tooling says so
--- \code{low\_fpr\_sample\_warning} exists for exactly this reason.
\item \textbf{Report unseen-generator results.} Seen-generator numbers flatter the
model and do not predict field behaviour.
\end{enumerate}

%==========================================================
\section{Open tickets}

\begin{center}
\begin{tabular}{@{}llp{7.4cm}@{}}
\toprule
\textbf{Ticket} & \textbf{Priority} & \textbf{Description}\\
\midrule
--- & P0 & Run research mode end to end; produce one real detection number\\
FUS-03 & P0 & Channel-aware reliability weighting (evidence is in hand)\\
--- & P0 & Fuse \code{cm\_score} into Track 1\\
DEMO-01/02 & P0 & Single-command startup, deterministic demo scenarios\\
BR-04 & P1 & Minimal prosody branch\\
CODEC-05 & P1 & Automated $\Delta$EER/$\Delta$AUC table (deltas are hand-computed)\\
CODEC-06 & P1 & Embedding preservation study; needs frozen encoders\\
UI-03/04 & P1 & Branch and channel explanations, live latency indicator\\
BR-05 & P2 & Identity branch adapter\\
\bottomrule
\end{tabular}
\end{center}

\begin{todobox}{Housekeeping before submission}
\begin{itemize}[nosep,leftmargin=1.4em]
\item README badges are stale: they read \code{tests-170\_passed} (now 183) and
\code{p95\_latency-6.7ms}. The 6.7\,ms figure is \emph{engine-only}; live transport
p95 is 82.53\,ms. A reviewer who notices will discount other claims.
\item \code{native/audio-runtime/target/} (Rust build artefacts) and
\code{data/manifests/} (absolute paths into \code{stork}) must be gitignored.
\item Two commits are unpushed and roughly 28 files uncommitted.
\end{itemize}
\end{todobox}

%==========================================================
\section{Honest summary for judging}

\textbf{Strong.} Streaming architecture with measured overload behaviour; channel
intelligence; codec robustness on a real corpus with unseen attacks; privacy by
construction; no mandatory paid APIs; every claim backed by a checked-in artefact.

\textbf{Weak.} The detector. Every metric is a placeholder at AUC $\approx 0.71$
clean and $0.53$ under noise. The repository states this everywhere rather than
obscuring it.

\begin{tcolorbox}[colback=strivegrey,colframe=strivenavy,
  title=\textbf{The single highest-leverage move}]
Get research mode running on \code{stork}'s RTX 3090. Corpus, manifests, codec
matrix, GPU and evaluation harness are all in place. What is missing is one model
bundle and one benchmark run --- and that converts ``a measurable streaming system''
into ``a measurable streaming system \emph{that detects voice clones}''.
\end{tcolorbox}

\vspace{0.8cm}
\begin{center}
\color{strivenavy}\itshape
Team STRIVE --- IISER Bhopal --- Smart India Hackathon 2026\\
\normalfont\small Secure Voices $\cdot$ Stronger India
\end{center}

\end{document}
